% sections/01-problem.tex --- Постановка наукової проблеми.

\citsection{Problem statement}

Language-model agents have advanced considerably in recent years. Multi-step reasoning, code generation, and tool invocation have all improved. The execution substrate on which agents run their actions, however, has not changed much. Most agent frameworks, following Toolformer and ReAct, run generated code in a single-threaded sandbox --- a Python interpreter, a container, or something equivalent. The model writes straight-line code, one process runs it, and the agent waits for as long as the script takes.

Single-threaded execution is suitable for retrieval, parsing, calculation, and light scripting. It falls short for workloads where the total work scales with sample count or problem size: Monte Carlo estimation, hyperparameter search, combinatorial optimisation, large-scale simulation, statistical bootstrap. Within a single process, such workloads either exhaust the time budget or require severe undersampling. In both cases, the bottleneck is computation, not the model's reasoning.

Providing the agent with a parallel runtime is the natural workaround. A problem arises, however: LLMs still struggle with parallel code. Benchmarks consistently show correctness dropping on MPI, OpenMP, and CUDA tasks relative to their serial equivalents, especially when memory access patterns become non-trivial. A model that handles serial code well often fails at the same algorithm when concurrency controls appear.

The two requirements --- parallel execution and synchronous-style code generation --- need not be mutually exclusive. A distributed runtime can transparently absorb parallelism, accepting per-worker logic written in straight-line code while assuming responsibility for fan-out, fan-in, and result aggregation. The PARCS framework constitutes such a runtime: a Kubernetes-based cluster of daemon workers that executes user-supplied modules at runtime. The remaining design problem is the agent-facing interface.

The interface considered in the present work is the Model Context Protocol (MCP), a recently introduced standard for agent--tool interaction. Tools advertised through MCP are usable by any compliant agent without bespoke integration. Positioning a distributed computing cluster behind MCP and exposing a compact set of tools aligned with the agent's natural stepwise decomposition allows agent-generated C\# code to retain a single-threaded coding paradigm, while parallel execution is transparently delivered by the platform. The research issue addressed in this paper is therefore the design of an interface between a language-model agent and a distributed parallel computing cluster that preserves the synchronous coding paradigm at the agent's surface while providing cluster-scale throughput beneath it.
